\documentclass[a4paper,twoside]{article}

\usepackage{morewrites}

\usepackage[svgnames,dvipsnames,x11names]{xcolor}
\usepackage{graphicx}
\usepackage[english]{babel}
\usepackage[colorlinks=true, linkcolor=NavyBlue, urlcolor=NavyBlue, citecolor=NavyBlue]{hyperref}
\usepackage[a4paper]{geometry}
\usepackage{ts-fonts}
\usepackage{float}
\usepackage{graphicx}
\usepackage{tcolorbox}
\usepackage{fvextra}
\usepackage{amsmath}
\usepackage{seqsplit}
\usepackage{ts-code}

\usepackage{lastpage}
\usepackage{refcount}
\makeatletter
\def\ps@plain{
\let\@oddhead\@empty
\let\@evenhead\@empty
\def\@oddfoot{
\hfil
\ifnum\getpagerefnumber{LastPage}>1\relax
\thepage
\fi
\hfil}
\let\@evenfoot\@oddfoot
}
\makeatother

\title{Mermaid Diagrams}
\author{}\date{}
\begin{document}
\maketitle
Just like MkDocs, TeXSmith understands \href{https://mermaid.js.org}{Mermaid} diagrams. Browsers render them on the fly, but PDFs need static assets, so TeXSmith converts each diagram into a vector PDF during the build. That requires either:
\begin{enumerate}
\item Installing \tscodeinline{mermaid-\allowbreak{}cli} and its dependencies in your environment, or
\item Using Docker with the pulled image \tscodeinline{mermaidjs/mermaid-\allowbreak{}cli}.
\end{enumerate}
\section{Inline diagram}
A Mermaid fence is a data directive whose node word is \tscodeinline{image}:
\begin{tscode}[lang=md, engine=pygments]
\begin{Verbatim}[commandchars=\\\{\},breaklines, breakanywhere, commandchars=\\\{\}]
\PY{l+s+sb}{```}\PY{l+s+sb}{mermaid}\PY{+w}{ }image
\PY{l+s}{flowchart LR}
\PY{l+s}{    A \PYZhy{}\PYZhy{}\PYZgt{} B}
\PY{l+s}{    B \PYZhy{}\PYZhy{}\PYZgt{} C}
\PY{l+s+sb}{```}
\end{Verbatim}
\end{tscode}
A bare \tscodeinline{```mermaid} fence means the same thing --- \tscodeinline{image} is the one
default that is not \tscodeinline{code}, because that is what MkDocs Material and every
Mermaid-aware renderer do with it. The sugar is kept indefinitely for exactly
that reason. To \emph{show} the source as a listing instead, ask for \tscodeinline{mermaid code}:
\begin{tscode}[lang=md, engine=pygments]
\begin{Verbatim}[commandchars=\\\{\},breaklines, breakanywhere, commandchars=\\\{\}]
\PY{l+s+sb}{```}\PY{l+s+sb}{mermaid}\PY{+w}{ }code
\PY{l+s}{flowchart LR}
\PY{l+s}{    A \PYZhy{}\PYZhy{}\PYZgt{} B}
\PY{l+s+sb}{```}
\end{Verbatim}
\end{tscode}
\begin{figure}[H]
\centering
\includegraphics[width=\linewidth]{<HASH>.pdf}
\end{figure}
\section{External diagrams}
Sometimes diagrams live better outside the Markdown. TeXSmith supports:
\begin{enumerate}
\item Reference external \tscodeinline{.mmd} / \tscodeinline{.mermaid} files.
\item Embed Mermaid Live snippets using \tscodeinline{pako:} URLs for live editing.
\end{enumerate}
Diagram sources \emph{are} images: the extension recognises the format from the
path and converts inline and external sources the same way.

Using a \tscodeinline{.mmd} file is as simple as including an image:
\begin{tscode}[lang=md, engine=pygments]
\begin{Verbatim}[commandchars=\\\{\},breaklines, breakanywhere, commandchars=\\\{\}]
![\PY{n+nt}{Build pipeline}](\PY{n+na}{../assets/mermaid.mmd})
\end{Verbatim}
\end{tscode}
\begin{figure}[H]
\centering
\includegraphics[width=\linewidth]{<HASH>.pdf}
\caption{Build pipeline}
\end{figure}
Mermaid Live encodes diagrams via Pako (a compression library) so you can share/edit them through URLs:
\begin{tscode}[lang=md, engine=pygments]
\begin{Verbatim}[commandchars=\\\{\},breaklines, breakanywhere, commandchars=\\\{\}]
![\PY{n+nt}{Online Diagram}](\PY{n+na}{https://mermaid.live/edit\PYZsh{}pako:eNpVTctugzAQ\PYZus{}BVrT4lEEMQEiA\PYZus{}tIemt7aE9tX}
\PY{n+na}{EODl4eSrAtY5q2iH8vEBGpe1jtzOzMdJBpicAgv\PYZhy{}hrVgrryPMbV2SYxg1o8T7uJVmtHoipsvNhXxWkNcfby8hMU}
\PY{n+na}{pV3O3E6iQKbx\PYZus{}6mVfmgcHjVHMaPxqE5vOgvJLm2V2El0Qon9jjXobnX\PYZus{}Iv4wGbO0GbxpOQSPChsJYE526IHNdpa}
\PY{n+na}{jBC60cjBlVgjBzacUtgzB676wWOE\PYZhy{}tS6nm1Wt0UJLBeXZkCtkcLhvhKFFfWdtagk2p1ulQO23tApBFgH38DCMPG}
\PY{n+na}{TMKI0TKKABkHswQ\PYZhy{}wlPrxOtrSKNnSOA2SsPfgd2oN\PYZus{}DTZ9H9\PYZus{}ZXFC})
\end{Verbatim}
\end{tscode}
When TeXSmith renders HTML/PDF it wraps the image with a link to the Mermaid Live editor. Click the preview to inspect the source:
\begin{figure}[H]
\centering
\includegraphics[width=\linewidth]{<HASH>.pdf}
\caption{Example Pako}
\end{figure}
\section{\LaTeX{} Rendering}
Here's how the diagrams look once TeXSmith embeds them:
\begin{figure}[H]
\centering
\includegraphics[width=\linewidth]{snippet-<HASH>.pdf}
\end{figure}
\section{Conversion by TeXSmith}
All Mermaid diagrams are converted to PDF and included with \tscodeinline{\textbackslash{}includegraphics}
so they integrate cleanly with templates and \LaTeX{} floats.

Add a caption line to promote the diagram to a numbered figure:
\begin{tscode}[lang=md, engine=pygments]
\begin{Verbatim}[commandchars=\\\{\},breaklines, breakanywhere, commandchars=\\\{\}]
![\PY{n+nt}{Build pipeline}](\PY{n+na}{../assets/mermaid.mmd})

Figure: The documentation build pipeline. \PYZob{}\PYZsh{}fig:pipeline\PYZcb{}
\end{Verbatim}
\end{tscode}
Printed output might deserve a different theme. Point \tscodeinline{mermaid\_config} to a JSON config (front matter or CLI \tscodeinline{-\allowbreak{}-\allowbreak{}attribute press.mermaid\_config=...}) to override:
\begin{tscode}[lang=yaml, engine=pygments]
\begin{Verbatim}[commandchars=\\\{\},breaklines, breakanywhere, commandchars=\\\{\}]
\PY{n+nn}{\PYZhy{}\PYZhy{}\PYZhy{}}
\PY{n+nt}{press}\PY{p}{:}
\PY{+w}{  }\PY{n+nt}{mermaid\PYZus{}config}\PY{p}{:}\PY{+w}{ }\PY{l+lScalar+lScalarPlain}{mermaid\PYZhy{}config.json}
\PY{n+nn}{\PYZhy{}\PYZhy{}\PYZhy{}}
\end{Verbatim}
\end{tscode}
Alternatively, you can add a \tscodeinline{mermaid-\allowbreak{}config.json} file to the \tscodeinline{\textasciitilde{}/.texsmith/} directory
to apply it globally to all your TeXSmith projects.
\end{document}
